% Chapter 3, Figure 10: the concept of free-stream velocity (scan: figures/ch3/fig10.png).
% The house model rocket (tamrfig `model', units of its length 6.4), nose left, with the velocity profiles above
% and below the body at the scan's station behind the nose joint. Computed by fig10.py: inside the boundary
% layer (laminar there) the profile is Blasius's, u = V f'(5y'/delta) from Table 1 (y' from the surface), with the
% thickness greatly exaggerated as the caption says (delta = 1.6 diameters, the scan's); V, the arrow rows and the
% station are the scan's. The boundary-layer edge below the body is a guide (dashed), as printed and as in
% Figure 13. The "Free stream" brace continues down into an arrow (the free stream extends outward), as printed.
% The axis is extended beyond the nose and the tail as a thin line, the centre line inside the body (as printed,
% and as Figures 1 and 12 and Chapter 1 Figures 6-8).
\documentclass[11pt]{standalone}
\usepackage{tamrfig}
% the profile family's styles (Ch3 Figs 6, 9, 10, 13, 16, 17, 18, 25): the computed profile u(y) as a data curve
% (s1, 1pt), its velocity arrows lighter with small heads; the base lines stay in ink
\tikzset{
  profile/.style={draw=s1, line width=1pt},
  profile arrow/.style={draw=s1, line width=0.5pt, -{Stealth[length=3.4pt, width=2.4pt]}},
  profile stroke/.style={draw=s1, line width=0.5pt},
  % a brace beside a profile, its tip pointing left at the label
  label brace/.style={draw=ink2, line width=0.5pt, decorate,
    decoration={brace, amplitude=4pt, raise=0pt}},
}
\makeatletter
% \csvline[<options>]{<file>}: the polyline through the x, y rows of a CSV file
\newcommand{\csvline}[2][]{%
  \pgfplotstableread[col sep=comma]{#2}\fig@tbl
  \pgfplotstablegetrowsof{\fig@tbl}\pgfmathtruncatemacro\fig@n{\pgfplotsretval-1}%
  \foreach \fig@i in {0,...,\fig@n}{%
    \pgfplotstablegetelem{\fig@i}{x}\of\fig@tbl\let\fig@x\pgfplotsretval
    \pgfplotstablegetelem{\fig@i}{y}\of\fig@tbl\let\fig@y\pgfplotsretval
    \ifnum\fig@i=0 \xdef\fig@path{(\fig@x,\fig@y)}\else\xdef\fig@path{\fig@path -- (\fig@x,\fig@y)}\fi}%
  \draw[#1] \fig@path;}
% \csvarrows[<options>]{<file>}: an arrow from (x0, y) to (x0 + u, y) for each row
\newcommand{\csvarrows}[2][]{%
  \pgfplotstableread[col sep=comma]{#2}\fig@tbl
  \pgfplotstablegetrowsof{\fig@tbl}\pgfmathtruncatemacro\fig@n{\pgfplotsretval-1}%
  \foreach \fig@i in {0,...,\fig@n}{%
    \pgfplotstablegetelem{\fig@i}{x0}\of\fig@tbl\let\fig@x\pgfplotsretval
    \pgfplotstablegetelem{\fig@i}{y}\of\fig@tbl\let\fig@y\pgfplotsretval
    \pgfplotstablegetelem{\fig@i}{u}\of\fig@tbl\let\fig@u\pgfplotsretval
    \draw[#1] (\fig@x,\fig@y) -- ++(\fig@u,0);}}
\makeatother
\begin{document}
\begin{tikzpicture}[x=1.5cm, y=1.5cm]
  \def\xs{2.33}\def\V{1.16}\def\R{0.2}\def\dl{0.65}\def\top{2.6}
  % ---- the rocket ----------------------------------------------------------------------------------------------
  \pic {rocket={model, centerline=false}};
  \draw[centerline] (0.1,0) -- (6.3,0);
  \draw[thin line] (-0.51,0) -- (0,0) (6.4,0) -- (6.91,0);       % the axis, extended (as Figures 1 and 12)
  % ---- the profiles: arrows, base lines, the lines through the arrow tips -----------------------------------
  \csvarrows[profile arrow]{figures/v2/ch3/fig10-arrows.csv}
  \draw[outline] (\xs,\R) -- (\xs,\top) (\xs,-\R) -- (\xs,-\top);
  \csvline[profile]{figures/v2/ch3/fig10-profile.csv}
  \csvline[profile, yscale=-1]{figures/v2/ch3/fig10-profile.csv}
  % V: the full length of a free-stream arrow, on the uppermost row
  \dimline{(\xs,2.435)}{(\xs+\V,2.435)}{$V$}
  % the edge of the boundary layer below the body
  \draw[guide] (\xs,-\R-\dl) -- (\xs+\V,-\R-\dl);
  % ---- the braces ------------------------------------------------------------------------------------------------
  \draw[label brace] (\xs-0.1,-\R-\dl+0.01) -- (\xs-0.1,-\R-0.01)
    node[midway, left=5pt] {Boundary layer};
  \draw[label brace] (\xs-0.1,-2.35) -- (\xs-0.1,-\R-\dl-0.02)
    node[midway, left=5pt] {Free stream};
  \draw[ink2, line width=0.5pt, -{Stealth[length=4pt,width=2.6pt]}] (\xs-0.1,-2.35) -- (\xs-0.1,-\top-0.05);
\end{tikzpicture}
\end{document}
